\section{Conclusions}
\label{sec:conclusions}

I have expanded the two-loop Type~II finite remainders
for $gg\to AA$ with two $O_G$ insertions about the
pair production threshold through $\bA^6$ at fixed
scattering angle. After the endpoint constants are
expressed in a basis of \BasisCountWord{}
transcendental constants, all inverse powers of
$\bA$ and the odd powers through $\bA^5$ cancel
identically. Every coefficient is even in
$\cos\theta$, as required by Bose symmetry.
For $0<\bA\leq\RecommendedBeta$, the maximum
relative errors of the two complete coefficients
are \PoneErrLePointTwo{} and
\PtwoErrLePointTwo{}, respectively, at the
points tested, including the forward angles.
The released series files are a factor of
\SizeRatio{} smaller than the expressions
released with Ref.~\cite{Bhattacharya:2019bmd}.
The timing comparisons in
App.~\ref{app:data} also show faster evaluation
for the representations tested.

A consistent comparison of the perturbative orders also required
converting the remainders of Ref.~\cite{Bhattacharya:2019bmd} to a common normalisation. The two-loop files carry an
overall factor of \FileFactor{} relative to
the tree and one-loop amplitudes, and the
one-loop files are given at $\muR=1$~GeV.
After the scale dependence is added and
all orders are brought to a common normalisation, the
remainders satisfy the renormalisation group
relations in Eqs.~\eqref{eq:rg-one} and
\eqref{eq:rg-two} within the accuracy of
the checks described in
App.~\ref{app:normalization}.
The Ref.~\cite{Bhattacharya:2019bmd} pole cancellation is unaffected, since all terms of the two-loop subtraction
carry the same overall factor.

At leading power, the angular dependence
enters entirely through the logarithm of
the transverse momentum of the pseudoscalars.
I identified the one-loop logarithm with
the DPS singularity studied in
Ref.~\cite{GauntStirling2011}.
The gluon splitting weights and the
relative signs of the CP-odd subprocess
amplitudes determine the momentum fraction
dependence of its coefficient.
With the overall constant fixed at
threshold, the resulting expression
reproduces the velocity dependence
of the exact one-loop coefficient
in the forward limit to a relative
accuracy of \DPSMaxDeviation{} over
the tested range.
The same analysis explains the absence
of a one-loop logarithm in the
$|J_z|=2$ amplitude, the independence
of the logarithmic coefficient from
$\nf$ and $\muR$, and its phase relative
to the tree amplitude.
An explicit one-loop calculation of
$gg\to HH$ through two CP-even vertices
confirms the corresponding scalar
weight $F_S$, using the normalisation
fixed from $AA$ and accounting for
the relative sign of the tree
amplitudes. The opposite helicity
amplitude again has no one-loop
DPS logarithm.

At two loops, the $J_z=0$ amplitude
grows as $\ln^3p_T$ and the
$|J_z|=2$ amplitude as $\ln^2p_T$.
The ratio of the leading two-loop
and one-loop logarithms in the
$J_z=0$ amplitude agrees with the
corresponding ratio in the all-orders
result for planar $\mathcal N=4$
super Yang--Mills
theory~\cite{DixonEsterlis2016}.
The coefficient of the quadratic
logarithm in this amplitude is
proportional to $\beta_0$, with
the structure discussed in
Sec.~\ref{sec:two-loop} in relation
to the NLO $g\to gg$ splitting
kernel for double parton
distributions~\cite{Diehl:2019rdh}.
The finite virtual corrections
are less singular than the
amplitudes. At leading power,
the one-loop DPS logarithm does
not contribute to $V_1$, because
it is out of phase with the
tree amplitude. The cubic
logarithm is absent from $V_2$,
whose highest logarithmic term is
\[
 V_2\supset
 -\frac74\pi^2\Nc^2
 \ln^2\!\left(\frac{2p_T}{\mA}\right).
\]
This coefficient results from a
partial cancellation between the
interference of the two-loop amplitude
with the tree amplitude and the square
of the one-loop amplitude.

Several questions remain open. The overall factor $\pi\Nc m^2$ multiplying the one-loop helicity weight has not yet been derived directly from the factorisation formula with the effective vertices. At two loops, a general treatment of the singular region would have to reproduce the different logarithmic powers of the $J_z=0$ and $|J_z|=2$ amplitudes, their colour and flavour dependence, and the reduction of the singularity in the finite virtual correction. In the forward limit, the two-loop logarithms of the $J_z=0$ amplitude follow the one-loop helicity weight at all velocities tested, and fits give numerical evidence that the real coefficient of the $|J_z|=2$ double logarithm is $-4\Nc$ times the one-loop forward amplitude, which itself has the closed form of Eq.~\eqref{eq:K1}; the mechanism producing this term has not been derived. The present results therefore provide explicit two-loop QCD benchmarks for such an analysis.
The known two-loop amplitudes for
$gg\to HH$ through two effective
vertices~\cite{Banerjee:2018lfq}
also permit a comparison with
the scalar case.
Finite top quark mass effects
are a separate question.
If the one-loop DPS factorisation
applies to the full theory, the
on-shell triangle form factors
in Eq.~\eqref{eq:top-dps} multiply
the logarithmic coefficient by
\TopMassDPSA{} at
$\mA=\TopMassMAA$~GeV.
This estimate does not determine
the mass dependence of the
remaining one-loop terms or the two-loop amplitude.

The series provide the two-loop Type~II amplitude input to the hard
function for the pair invariant mass distribution near $m_{AA}=2\mA$.
They enter at NNLO when the Type~II tree amplitude is the leading
amplitude, as for a pseudoscalar coupled to gluons only through
$G\tilde G$, and at N$^3$LO in a 2HDM, through interference with the
leading Type~I amplitude; Sec.~\ref{sec:use} describes how they are
combined with the other contributions. The mixed $O_G$--$O_J$
contribution, omitted from the numerical results for $V_2$, gives a
kinematically constant shift and leaves the logarithmic coefficients
unchanged. The series, the evaluator, the validation records and the code
of the construction are provided with the paper~\cite{SupplementalMaterial}.
